% ============================================================================
%  D/16-gan-nghia-va-van-vuc — file chương
%  Ngày tạo: 2026-09-07 | Sửa gần nhất: 2026-09-07 | Ôn gần nhất: chưa
%  Dependency: A/01, A/02, A/04. Mức độ: cốt lõi.
% ============================================================================
\documentclass[../../english.tex]{subfiles}
\begin{document}
\chapter{Gần nghĩa và văn vực (Near-Synonyms and Register)}
\ptrtarget{D/16}\ptrtarget{D/16-gan-nghia-va-van-vuc}
\dependency{\ptr{A/01} cho ý ``nghĩa nằm ở cách dùng''; \ptr{A/02} cho chín mặt của việc biết một từ; \ptr{A/04} cho nguồn gốc từ.}

\begin{tomtat}
Chương này xử lý câu hỏi thực dụng nhất của người học ở trình độ trung cấp trở lên: từ điển cho ba từ tiếng Anh cho một từ tiếng Việt, chọn cái nào? Ý trung tâm: hai từ ``đồng nghĩa'' hầu như không bao giờ thay thế được nhau hoàn toàn, và cái tách chúng thường không phải nghĩa mà là \emph{văn vực}\cite{biber1999} --- mức trang trọng, ngành nào dùng, và các từ mà nó quen đi cùng. Nội dung gồm: bốn chiều tách các từ gần nghĩa; vì sao tiếng Anh có nhiều cặp gần nghĩa đến vậy (lý do lịch sử, và nó dự đoán được sắc thái); màu sắc đánh giá ẩn trong từ; và cách tra cứu bằng dữ liệu ngôn ngữ thật thay vì đoán. Nếu chỉ nhớ một điều: đừng hỏi ``từ này nghĩa là gì'' mà hỏi ``ai dùng từ này, trong hoàn cảnh nào, và nó thường đi với từ nào''.
\end{tomtat}

\begin{nentang}
\kn{gần nghĩa}{hai từ có nghĩa cốt lõi gần nhau nhưng khác về văn vực, sắc thái, hoặc phạm vi dùng.}
\kn{văn vực}{mức trang trọng và bối cảnh dùng: nói thường, viết chung, học thuật, kỹ thuật, pháp lý.}
\kn{màu sắc đánh giá}{thái độ tích cực hoặc tiêu cực gắn liền với từ, dù nghĩa mô tả trung tính.}
\kn{phạm vi}{loại đối tượng mà từ áp dụng được: \emph{high} đi với \emph{temperature} nhưng không với \emph{a tall building}.}
\kn{cặp Anh--Latin}{hiện tượng tiếng Anh có một từ gốc Germanic và một từ gốc Latin/Pháp cho cùng khái niệm.}
\kn{ngữ liệu}{kho văn bản thật dùng để tra xem một từ thực sự được dùng thế nào và với những từ nào.}
\end{nentang}

\begin{khoigoi}
\h{\emph{Begin}, \emph{start}, \emph{commence}, \emph{initiate} --- cái gì tách chúng, nếu không phải nghĩa?}
\h{Vì sao tiếng Anh lại có nhiều cặp gần nghĩa hơn hầu hết các ngôn ngữ châu Âu khác?}
\h{Hai từ có nghĩa mô tả giống hệt nhau mà một từ khen một từ chê --- cơ chế nào tạo ra chuyện đó?}
\h{Bạn không chắc chọn \emph{significant} hay \emph{considerable}. Có cách nào kiểm tra ngoài việc hỏi người bản ngữ không?}
\h{Vì sao từ điển song ngữ lại đặc biệt yếu ở chỗ này?}
\end{khoigoi}

Hãy bắt đầu bằng bốn từ ở câu hỏi mở đầu thứ nhất. Từ điển sẽ dịch cả bốn là ``bắt đầu''. Nhưng thử đặt chúng vào câu:

\begin{tuvung}[Bốn cách nói ``bắt đầu'']
\tv{start}{Bắt đầu}{Phổ thông nhất, dùng được gần như mọi chỗ; cũng dùng cho máy móc}{Let's start. / The engine started.}
\tv{begin}{Bắt đầu}{Hơi trang trọng hơn \emph{start}; thiên về quá trình hơn là hành động bật lên}{The lecture begins at nine.}
\tv{commence}{Bắt đầu}{Rất trang trọng, gần văn bản pháp lý và nghi lễ; nghe cứng trong văn nói}{Proceedings will commence at noon.}
\tv{initiate}{Khởi động}{Kỹ thuật và hành chính; ngụ ý có người hoặc hệ thống chủ động khởi phát}{The handshake is initiated by the client.}
\end{tuvung}

Không cặp nào khác nhau về nghĩa cốt lõi. Cái tách chúng là ba thứ khác: mức trang trọng, ngành nào dùng, và loại chủ ngữ đi được. Đó là nội dung của cả chương.

\section{Bốn chiều tách các từ gần nghĩa}

\emph{Chiều một --- mức trang trọng.} Đây là chiều dễ thấy nhất và cũng là chiều gây lỗi nhiều nhất, vì người học có xu hướng chọn từ trang trọng để nghe ``giỏi''. Kết quả là văn nghe cứng và đôi khi buồn cười: dùng \emph{commence} trong một email nội bộ.

\emph{Chiều hai --- ngành và loại văn bản.} \emph{Initiate} thuộc về kỹ thuật và hành chính; \emph{demonstrate} trong khoa học có nghĩa hẹp hơn trong đời thường. Một từ có thể hoàn toàn tự nhiên trong ngành này và lạc lõng trong ngành khác.

\emph{Chiều ba --- phạm vi áp dụng.} Đây là chiều mà người học ít để ý nhất. \emph{High} và \emph{tall} đều là ``cao'', nhưng \emph{high} đi với \emph{temperature}, \emph{pressure}, \emph{cost}, còn \emph{tall} đi với người và toà nhà. Không có logic nào ở đây --- chỉ có thói quen, và đó là cầu sang \ptr{D/17}.

\emph{Chiều bốn --- màu sắc đánh giá.} \emph{Ambitious} và \emph{unrealistic} có thể mô tả cùng một kế hoạch. \emph{Careful} và \emph{timid} có thể mô tả cùng một người. Chiều này quan trọng đặc biệt khi bạn viết về công trình của người khác (\ptr{C/13} mục 4).

\begin{figure}[H]\centering
\begin{tikzpicture}[
  ax/.style={-{Latex[length=2.4mm]},draw=steel,very thick},
  w/.style={draw=steel,fill=bluebg,rounded corners=2pt,inner sep=3pt,font=\footnotesize},
  lb/.style={font=\scriptsize\itshape,color=tiercol,align=center}]
\draw[ax] (0,0) -- (11.0,0);
\node[lb,below=0.5cm] at (0.9,0) {nói thường};
\node[lb,below=0.5cm] at (4.2,0) {viết chung};
\node[lb,below=0.5cm] at (7.4,0) {học thuật};
\node[lb,below=0.5cm] at (10.2,0) {pháp lý / nghi lễ};
\node[w] at (0.9,0.45) {start};
\node[w] at (3.4,0.45) {begin};
\node[w] at (6.0,0.45) {initiate};
\node[w] at (8.6,0.45) {commence};
\node[w] at (0.9,1.15) {get};
\node[w] at (3.4,1.15) {obtain};
\node[w] at (6.0,1.15) {acquire};
\node[w] at (8.6,1.15) {procure};
\node[w] at (0.9,1.85) {ask};
\node[w] at (3.4,1.85) {request};
\node[w] at (6.0,1.85) {enquire};
\node[w] at (8.6,1.85) {petition};
\end{tikzpicture}
\caption{Ba bộ từ gần nghĩa xếp theo chiều trang trọng. Để ý: cột trái gần như toàn từ gốc Germanic ngắn, cột phải toàn từ gốc Latin dài --- không phải trùng hợp, xem mục 3.}
\label{fig:van-vuc-thang}
\end{figure}

\section{Màu sắc đánh giá}

Câu hỏi mở đầu thứ ba hỏi cơ chế tạo ra hai từ cùng nghĩa mô tả mà khác thái độ. Cơ chế là \emph{lây nhiễm từ ngữ cảnh}: một từ liên tục xuất hiện cạnh những từ tiêu cực sẽ dần mang màu tiêu cực, kể cả khi định nghĩa của nó trung tính.

Hiện tượng này được mô tả trong ngôn ngữ học ngữ liệu dưới tên \emph{semantic prosody}\cite{louw1993}. Ví dụ kinh điển là \emph{cause}. Nghĩa từ điển hoàn toàn trung tính, nhưng trong sử dụng thực tế nó đi với \emph{problems}, \emph{damage}, \emph{delay}, \emph{failure} nhiều hơn hẳn với những thứ tốt --- nên nó mang màu tiêu cực nhẹ. So sánh với \emph{lead to}, trung tính hơn, và \emph{bring about}, thường dùng cho thay đổi được coi là tích cực.

Điều quan trọng: bạn không tra được màu sắc này trong từ điển vì nó không nằm trong định nghĩa. Nó chỉ hiện ra khi bạn nhìn nhiều ví dụ thật --- đó là mục 5.

\begin{tuvung}[Cùng nghĩa mô tả, khác thái độ]
\tv{ambitious}{Tham vọng}{Khen: mục tiêu cao và đáng theo đuổi}{An ambitious research programme.}
\tv{unrealistic}{Không thực tế}{Chê: cùng mục tiêu đó, nhưng bị coi là không khả thi}{An unrealistic research programme.}
\tv{thorough}{Kỹ lưỡng}{Khen: xét đủ mọi mặt}{A thorough evaluation.}
\tv{exhaustive}{Vét cạn}{Trung tính đến hơi nặng nề; nhấn khối lượng hơn chất lượng}{An exhaustive list of parameters.}
\tv{simple}{Đơn giản}{Khen trong kỹ thuật: ít bộ phận, dễ hiểu}{A simple and effective method.}
\tv{simplistic}{Giản lược quá mức}{Chê: bỏ qua cái quan trọng}{A simplistic model of the noise.}
\tv{robust}{Bền vững}{Khen: chịu được điều kiện xấu}{Robust to illumination changes.}
\tv{rigid}{Cứng nhắc}{Chê: không thích ứng được}{A rigid pipeline that fails on new data.}
\end{tuvung}

\section{Gốc rễ: vì sao tiếng Anh nhiều cặp gần nghĩa}

Câu hỏi mở đầu thứ hai có một câu trả lời lịch sử, và nó không chỉ là chuyện thú vị --- nó cho bạn một quy tắc dự đoán dùng được ngay.

Tiếng Anh nhận hai đợt từ vựng ngoại lai lớn. Đợt thứ nhất bắt đầu sau năm 1066, khi tiếng Pháp Norman trở thành ngôn ngữ của triều đình, luật pháp và tầng lớp cai trị ở Anh trong vài thế kỷ. Tiếng Anh bản địa --- gốc Germanic --- vẫn là ngôn ngữ của đời sống hằng ngày. Kết quả là hàng loạt cặp từ hình thành, trong đó từ Germanic dùng cho việc thường ngày và từ Pháp dùng cho việc trang trọng hoặc thuộc quyền lực. Ví dụ hay được nhắc: con vật nuôi mang tên Germanic (\emph{cow}, \emph{pig}, \emph{sheep}) còn món ăn trên bàn tầng lớp trên mang tên Pháp (\emph{beef}, \emph{pork}, \emph{mutton}).

Đợt thứ hai đến trong thời Phục hưng, khi giới học giả mượn ồ ạt từ tiếng Latin và Hy Lạp để diễn đạt các khái niệm học thuật. Đợt này tạo ra tầng từ vựng học thuật mà \ptr{A/02} đã nhắc, và nó chồng lên tầng Pháp sẵn có.

Kết quả là tiếng Anh thường có \emph{ba} tầng cho cùng một khái niệm: Germanic ngắn cho đời thường, Pháp cho mức trung, Latin dài cho học thuật. \emph{Ask --- question --- interrogate}. \emph{Rise --- mount --- ascend}. \emph{Fire --- flame --- conflagration}.

Từ đó ra quy tắc dự đoán rất hữu ích: \textbf{từ càng dài và càng có gốc Latin thì càng trang trọng và càng hẹp phạm vi}. Không phải luật tuyệt đối, nhưng đúng đủ thường xuyên để dùng khi bạn chưa tra được. Và nó nối thẳng với \ptr{A/04}: nhận ra gốc từ không chỉ giúp đoán nghĩa mà còn giúp đoán \emph{văn vực}, và nguồn tra gốc từ tiện nhất là \cite{etymonline}.

Một hệ quả thực hành đáng nhớ: các động từ hai phần kiểu \emph{put off}, \emph{carry out}, \emph{look into} phần lớn có gốc Germanic và thuộc văn vực nói hoặc viết chung; mỗi cái thường có một từ Latin tương ứng thuộc văn vực học thuật --- \emph{postpone}, \emph{conduct}, \emph{investigate}. Đây là lý do văn học thuật tiếng Anh dùng ít động từ hai phần, và cũng là một cách nhanh để nâng văn vực khi cần.

\begin{gocre}
\gr{Sau 1066 --- tiếng Pháp Norman}{Ai nói ngôn ngữ nào ở Anh?}{Tầng Germanic cho đời thường, tầng Pháp cho quyền lực và trang trọng.}
\gr{Phục hưng --- mượn Latin và Hy Lạp\cite{etymonline}}{Diễn đạt khái niệm học thuật mới thế nào?}{Tầng thứ ba, dài và hẹp \(\Rightarrow\) từ vựng học thuật ở \ptr{A/02}.}
\gr{Ba tầng chồng nhau}{Vì sao nhiều từ gần nghĩa đến vậy?}{Quy tắc dự đoán: dài \(+\) gốc Latin \(=\) trang trọng hơn, hẹp hơn.}
\gr{Động từ hai phần}{Vì sao văn học thuật ít dùng chúng?}{Chúng thuộc tầng Germanic; mỗi cái có bản Latin tương ứng thuộc tầng học thuật.}
\gr{Từ điển song ngữ}{Vì sao chúng không tách được các từ này?}{Chúng ánh xạ nghĩa, không ánh xạ văn vực --- xem mục 6.}
\gr{Ngữ liệu ngôn ngữ\cite{coca}}{Làm sao biết một từ thật sự được dùng thế nào?}{Đếm trên kho văn bản thật; hiện ra cả văn vực lẫn màu sắc đánh giá.}
\end{gocre}

\section{Cách tra bằng dữ liệu thật}

Câu hỏi mở đầu thứ tư hỏi cách kiểm tra mà không cần người bản ngữ. Có, và nó dựa trên đúng ý của \ptr{A/01}: nghĩa nằm ở cách dùng, nên muốn biết cách dùng thì nhìn \emph{nhiều} lần dùng thật.

\emph{Cách một --- từ điển học thuật cho người học}\cite{oald,ldoce}\emph{.} Khác với từ điển song ngữ, loại này ghi nhãn văn vực (\emph{formal}, \emph{technical}, \emph{informal}), cho ví dụ câu thật, và liệt kê các từ hay đi cùng. Đây nên là từ điển mặc định của bạn.

\emph{Cách hai --- công cụ ngữ liệu}\cite{coca}\emph{.} Chúng cho bạn tra một từ và xem hàng trăm câu thật chứa nó, kèm thống kê những từ hay đứng cạnh. Đây là cách duy nhất phát hiện được màu sắc đánh giá ở mục 2: nếu \emph{cause} liên tục đứng cạnh \emph{problem}, \emph{damage}, \emph{delay}, bạn thấy ngay màu của nó.

\emph{Cách ba --- ngữ liệu của chính ngành bạn.} Đây là cách rẻ nhất và chính xác nhất cho mục đích của bạn. Tải về ba mươi bài báo trong lĩnh vực của bạn, gộp thành một tệp văn bản, rồi tìm kiếm từ bạn phân vân. Bạn sẽ thấy ngay từ nào được dùng, trong cấu trúc nào, và với những từ nào --- đúng cộng đồng bạn viết cho, chứ không phải tiếng Anh chung.

\begin{cannho}
Phép thử ba câu: khi phân vân giữa hai từ, tìm ba câu thật chứa mỗi từ trong tài liệu ngành của bạn. Nếu một từ không xuất hiện lần nào trong ba mươi bài, đừng dùng nó --- dù từ điển nói nó đúng nghĩa. Đây là cách kiểm tra nhanh nhất và đáng tin nhất mà không cần hỏi ai.
\end{cannho}

\section{Vì sao từ điển song ngữ yếu ở đây}

Câu hỏi mở đầu thứ năm. Ba lý do, và hiểu chúng giúp bạn dùng từ điển song ngữ đúng chỗ thay vì bỏ hẳn.

\emph{Lý do một --- nó ánh xạ nghĩa, không ánh xạ văn vực.} Cả bốn từ ở mục đầu đều dịch là ``bắt đầu'', và điều đó không sai --- nó chỉ không đủ.

\emph{Lý do hai --- ranh giới nghĩa giữa hai ngôn ngữ không trùng nhau.} Một từ tiếng Việt có thể phủ vùng nghĩa mà tiếng Anh chia cho ba từ, và ngược lại. Ánh xạ một--nhiều này không phải khiếm khuyết của từ điển mà là bản chất của việc dịch.

\emph{Lý do ba --- nó không cho biết từ nào đi với từ nào.} Và như \ptr{D/17} sẽ chỉ ra, đây thường là thông tin quyết định.

Kết luận thực hành không phải ``bỏ từ điển song ngữ'' mà là: dùng nó để \emph{tìm ứng viên}, rồi dùng từ điển cho người học hoặc ngữ liệu để \emph{chọn} trong số ứng viên đó. Hai bước, hai công cụ.

\section{Liên hệ ngang}

\begin{lienhe}
\lh{Lịch sử}{Cuộc chinh phục Norman}{Giải thích trực tiếp cấu trúc ba tầng của từ vựng tiếng Anh; một sự kiện chính trị để lại dấu vết trong mọi câu bạn viết hôm nay.}
\lh{Xã hội học}{Ngôn ngữ và địa vị}{Chọn từ trang trọng là một tín hiệu xã hội, không chỉ là chọn nghĩa --- lý do dùng sai văn vực gây phản ứng mạnh hơn sai ngữ pháp.}
\lh{Học máy}{Biểu diễn từ theo ngữ cảnh}{Cùng giả định với \ptr{A/01}: nghĩa của từ suy ra được từ các từ xung quanh nó. Các mô hình ngôn ngữ học sắc thái theo đúng cách mục 4 mô tả.}
\lh{Tiếng Việt}{Từ Hán Việt và từ thuần Việt}{Cấu trúc song song đáng kinh ngạc: ``bắt đầu''/``khởi đầu'', ``hỏi''/``chất vấn''. Cùng quy tắc: gốc vay mượn trang trọng hơn.}
\lh{Luật}{Thuật ngữ pháp lý}{Nhiều từ pháp lý tiếng Anh giữ nguyên tầng Pháp Norman --- di sản trực tiếp của mục 3, và lý do văn bản pháp lý nghe khác hẳn văn thường.}
\lh{Marketing}{Chọn từ để định khung}{Khai thác có chủ đích màu sắc đánh giá ở mục 2; đọc ra kỹ thuật này là một phần của \ptr{B/08}.}
\end{lienhe}

\begin{traloi}
\tl{Không phải nghĩa cốt lõi --- cả bốn đều là ``bắt đầu''. Cái tách chúng là ba thứ khác: mức trang trọng (\emph{start} phổ thông, \emph{commence} gần văn bản pháp lý), ngành và loại văn bản (\emph{initiate} thuộc kỹ thuật và hành chính), và loại chủ ngữ đi được. Nói tổng quát, các từ gần nghĩa được tách bởi bốn chiều: trang trọng, ngành, phạm vi áp dụng, và màu sắc đánh giá.}
\tl{Vì tiếng Anh nhận hai đợt vay mượn lớn chồng lên tầng bản địa. Sau 1066, tiếng Pháp Norman trở thành ngôn ngữ triều đình và luật pháp trong vài thế kỷ trong khi tiếng Anh Germanic giữ vai trò đời thường. Rồi thời Phục hưng, giới học giả mượn ồ ạt từ Latin và Hy Lạp. Kết quả là ba tầng cho cùng khái niệm: \emph{ask --- question --- interrogate}. Quy tắc dự đoán rút ra: từ càng dài, càng gốc Latin thì càng trang trọng và càng hẹp phạm vi.}
\tl{Lây nhiễm từ ngữ cảnh: một từ liên tục xuất hiện cạnh những từ tiêu cực dần mang màu tiêu cực, dù định nghĩa của nó trung tính. \emph{Cause} là ví dụ kinh điển --- nó đi với \emph{problems}, \emph{damage}, \emph{delay} nhiều hơn hẳn với những thứ tốt, nên mang màu tiêu cực nhẹ so với \emph{lead to}. Bạn không tra được màu này trong từ điển vì nó không nằm trong định nghĩa; nó chỉ hiện ra khi nhìn nhiều ví dụ thật.}
\tl{Có ba cách. Từ điển học thuật cho người học --- nó ghi nhãn văn vực, cho ví dụ thật, và liệt kê từ hay đi cùng. Công cụ ngữ liệu --- cho bạn xem hàng trăm câu thật kèm thống kê từ đứng cạnh, và đây là cách duy nhất phát hiện màu sắc đánh giá. Và rẻ nhất cho mục đích của bạn: gộp ba mươi bài báo trong ngành thành một tệp rồi tìm kiếm từ đó. Phép thử ba câu: nếu từ không xuất hiện lần nào trong ba mươi bài, đừng dùng.}
\tl{Ba lý do. Nó ánh xạ nghĩa chứ không ánh xạ văn vực --- cả bốn từ ``bắt đầu'' đều dịch đúng mà vẫn không đủ. Ranh giới nghĩa giữa hai ngôn ngữ không trùng nhau, nên ánh xạ luôn là một--nhiều. Và nó không cho biết từ nào đi với từ nào, thứ thường quyết định (\ptr{D/17}). Kết luận không phải bỏ nó mà là dùng nó để \emph{tìm ứng viên}, rồi dùng từ điển cho người học hoặc ngữ liệu để \emph{chọn}.}
\end{traloi}

\begin{dongian}
Có một chỗ mà từ điển song ngữ sẽ làm bạn thất vọng, và biết trước thì đỡ mất thời gian: khi nó cho bạn ba, bốn từ tiếng Anh cho cùng một từ tiếng Việt.

Ví dụ ``bắt đầu'' cho ra \emph{start}, \emph{begin}, \emph{commence}, \emph{initiate}. Cả bốn đều đúng nghĩa. Nhưng \emph{start} là từ bạn nói với bạn bè, \emph{commence} là từ trong văn bản pháp lý, còn \emph{initiate} là từ trong tài liệu kỹ thuật. Dùng \emph{commence} trong email nội bộ thì nghe kỳ, dù không sai chữ nào.

Nên câu hỏi đúng không phải ``từ này nghĩa là gì'' mà là ``ai dùng từ này, ở đâu''.

Có một quy tắc đoán nhanh, và nó đến từ lịch sử. Tiếng Anh có ba tầng từ vựng: tầng bản địa ngắn gọn (\emph{ask}), tầng gốc Pháp từ thời Norman (\emph{question}), và tầng gốc Latin mượn thời Phục hưng (\emph{interrogate}). Quy tắc: \emph{từ càng dài, gốc Latin càng rõ, thì càng trang trọng và càng hẹp}. Không tuyệt đối, nhưng đúng đủ nhiều để dùng khi bí.

Một chuyện nữa ít ai nói: nhiều từ mang sẵn thái độ khen hay chê, dù nghĩa mô tả trung tính. \emph{Ambitious} và \emph{unrealistic} có thể tả cùng một kế hoạch --- một từ khen, một từ chê. \emph{Simple} khen, \emph{simplistic} chê. Từ điển không ghi những chuyện này.

Vậy kiểm tra thế nào mà không cần hỏi người bản ngữ? Có một cách rẻ và rất hiệu quả: tải ba mươi bài báo trong đúng ngành bạn, gộp lại thành một tệp, rồi tìm kiếm từ bạn đang phân vân. Bạn sẽ thấy ngay từ nào được dùng, trong câu như thế nào. Và nếu một từ không xuất hiện lần nào trong ba mươi bài đó --- đừng dùng nó, dù từ điển nói nó đúng.
\motcau{Đừng hỏi từ này nghĩa gì; hỏi ai dùng nó ở đâu.}
\chuadongian{Quy tắc ``dài và Latin thì trang trọng hơn'' có nhiều ngoại lệ, và biết ngoại lệ nào chỉ đến từ tiếp xúc nhiều với ngành.}
\end{dongian}

\begin{onbai}
\ar[3]{Bốn chiều tách từ gần nghĩa}{Mức trang trọng; ngành và loại văn bản; phạm vi áp dụng; màu sắc đánh giá.}
\ar[3]{Câu hỏi đúng về một từ}{Không phải ``nghĩa gì'' mà ``ai dùng, ở đâu, đi với từ nào''.}
\ar[3]{Ba tầng từ vựng tiếng Anh}{Germanic bản địa; Pháp Norman sau 1066; Latin/Hy Lạp thời Phục hưng.}
\ar[3]{Quy tắc dự đoán văn vực}{Càng dài, càng gốc Latin thì càng trang trọng và càng hẹp phạm vi.}
\ar[2]{Ví dụ ba tầng}{\emph{ask --- question --- interrogate}; \emph{rise --- mount --- ascend}.}
\ar[2]{Động từ hai phần}{Gốc Germanic, thuộc văn nói/viết chung; mỗi cái có bản Latin học thuật (\emph{carry out}/\emph{conduct}).}
\ar[2]{Màu sắc đánh giá}{Sinh ra do lây nhiễm ngữ cảnh; \emph{cause} hơi tiêu cực vì đi với \emph{problem}, \emph{damage}.}
\ar[2]{Cặp khen--chê}{\emph{ambitious}/\emph{unrealistic}; \emph{simple}/\emph{simplistic}; \emph{robust}/\emph{rigid}.}
\ar[2]{Ba cách tra bằng dữ liệu thật}{Từ điển cho người học; công cụ ngữ liệu; và gộp ba mươi bài trong ngành bạn.}
\ar[2]{Phép thử ba câu}{Tìm ba câu thật chứa từ đó trong tài liệu ngành; không có lần nào thì đừng dùng.}
\ar[2]{Ba yếu điểm của từ điển song ngữ}{Không ánh xạ văn vực; ranh giới nghĩa hai ngôn ngữ không trùng; không cho biết từ đi với từ nào.}
\ar[1]{Cách dùng từ điển song ngữ đúng}{Dùng nó tìm ứng viên, rồi dùng công cụ khác để chọn.}
\end{onbai}

\begin{hoidap}
\qa{Tôi có nên học thuộc các bảng từ đồng nghĩa không?}{Không, vì bảng đồng nghĩa cho bạn danh sách mà không cho bạn tiêu chí chọn --- đúng cái vấn đề mà chương này nói. Cách hiệu quả hơn là học \emph{theo cặp có ngữ cảnh}: mỗi khi gặp một từ gần nghĩa với từ bạn đã biết, ghi cả hai vào cùng một mục trong bảng thuật ngữ (\ptr{B/10}) kèm một câu thật cho mỗi từ. Cặp có ngữ cảnh nhớ được và dùng được; danh sách trần thì không.}
\qa{Dùng từ ``sang'' có làm bài của tôi nghe chuyên nghiệp hơn không?}{Thường là ngược lại. Người đọc bản ngữ có kinh nghiệm nhận ra ngay khi một từ trang trọng bị đặt sai văn vực, và hiệu ứng không phải ``người này giỏi tiếng Anh'' mà là ``người này không quen với loại văn bản này''. Ngoài ra nó xung đột với \ptr{C/11}: từ Latin dài thường kéo theo danh từ hoá và câu nặng. Tiêu chí đúng không phải sang hay không mà là \emph{có xuất hiện trong tài liệu ngành bạn không}.}
\qa{Tiếng Anh Anh và tiếng Anh Mỹ khác nhau nhiều không ở chỗ này?}{Có khác, nhưng ít hơn bạn lo với văn học thuật. Khác biệt chủ yếu ở chính tả, một số từ vựng đời thường, và mức ưa dùng một vài cấu trúc. Với viết khoa học, quy tắc thực dụng: chọn một chuẩn và nhất quán trong cả bài, và nếu nơi bạn gửi bài có quy định thì theo quy định đó. Đừng dành nhiều công cho câu hỏi này --- nó ít ảnh hưởng tới việc người đọc hiểu bạn.}
\qa{Mô hình ngôn ngữ có giúp chọn từ được không?}{Có, và đây là một trong những việc chúng làm tốt --- vì chúng học chính xác cái mà mục 4 mô tả: từ nào đi với từ nào trong bao nhiêu ngữ cảnh. Cách hỏi hiệu quả là hỏi \emph{so sánh}: ``trong văn viết kỹ thuật, khác biệt giữa X và Y là gì, cho ví dụ'' --- thay vì ``từ nào đúng''. Nhưng vẫn kiểm lại bằng phép thử ba câu trong tài liệu ngành bạn, vì mô hình nói về tiếng Anh chung chứ không về cộng đồng cụ thể của bạn. Xem thêm \ptr{E/20}.}
\qa{Làm sao tôi biết một từ mang màu tiêu cực?}{Không có cách nào ngoài việc nhìn nhiều ví dụ, vì thông tin này không nằm trong định nghĩa. Cách nhanh: tra từ đó trong công cụ ngữ liệu và đọc lướt hai chục câu, chú ý các danh từ đi sau nó. Nếu chúng nghiêng hẳn về một phía --- toàn thứ xấu hoặc toàn thứ tốt --- bạn đã có câu trả lời. Với người học ở trình độ của bạn, đây là một trong những lý do mạnh nhất để dùng ngữ liệu thay vì chỉ dùng từ điển.}
\qa{Tiếng Việt cũng có hiện tượng ba tầng này không?}{Có một cấu trúc song song rất rõ: từ thuần Việt so với từ Hán Việt. ``Bắt đầu'' và ``khởi đầu'', ``hỏi'' và ``chất vấn'', ``nước'' và ``quốc gia''. Quy tắc cũng giống: lớp vay mượn trang trọng hơn, hẹp hơn, và hay dùng trong văn viết. Nhận ra sự song song này có ích thực tế: nó cho bạn một trực giác sẵn có để chuyển sang tiếng Anh, thay vì phải xây từ đầu.}
\qa{Tôi nên ưu tiên học chiều nào trong bốn chiều?}{Chiều ba --- phạm vi áp dụng --- nếu mục tiêu là viết đúng, vì lỗi ở chiều này gây câu nghe sai rõ nhất và người đọc bản ngữ phản ứng mạnh nhất. Chiều bốn --- màu sắc đánh giá --- nếu mục tiêu là đọc hiểu đúng ý người viết, vì đó là chỗ thái độ của tác giả lộ ra (\ptr{B/08}). Chiều một dễ tự học qua tiếp xúc; chiều hai bạn sẽ tự có khi xây bảng thuật ngữ ngành (\ptr{B/10}).}
\end{hoidap}

\begin{baitap}
\bt[1]{Xếp 4 bộ từ gần nghĩa theo thang trang trọng như hình ở mục 1}{Tự chọn}{Kiểm lại bằng từ điển cho người học có ghi nhãn văn vực.}
\bt[1]{Tìm 5 cặp khen--chê cùng nghĩa mô tả trong ngành bạn}{Tài liệu ngành}{Ví dụ mẫu ở mục 2; tìm cặp riêng của ngành bạn.}
\bt[2]{Gộp 30 bài báo ngành bạn thành một tệp và tra 10 từ bạn hay phân vân}{Tài liệu ngành}{Ghi lại từ nào không xuất hiện lần nào.}
\bt[2]{Với 5 động từ hai phần bạn hay dùng, tìm bản Latin tương ứng}{Bài viết của bạn}{\emph{carry out}/\emph{conduct}, \emph{look into}/\emph{investigate}\ldots}
\bt[2]{Tra \emph{cause}, \emph{lead to}, \emph{bring about} trong ngữ liệu và ghi 10 danh từ đi sau mỗi từ}{Công cụ ngữ liệu}{So ba danh sách để thấy màu sắc đánh giá.}
\bt[3]{Lập bảng 20 cặp gần nghĩa của ngành bạn, mỗi từ kèm một câu thật}{Tài liệu ngành}{Gộp vào bảng thuật ngữ ở \ptr{B/10}, trường ``từ lân cận''.}
\bt[3]{Lấy một đoạn bạn viết và thay 10 từ bằng từ gần nghĩa, rồi nhờ người đọc so hai bản}{Bài viết của bạn}{Ghi lại chỗ nào người đọc thấy khác biệt --- đó là chỗ sắc thái thật sự quan trọng.}
\end{baitap}

\section*{Kết chương và đọc tiếp}

Chương này cho bốn chiều tách các từ gần nghĩa, một quy tắc lịch sử để đoán văn vực khi chưa tra được, và một quy trình kiểm bằng dữ liệu thật thay vì đoán.

Nhưng ở chiều ba --- phạm vi áp dụng --- ta đã chạm vào một hiện tượng chưa giải thích được: vì sao \emph{high temperature} tự nhiên còn \emph{tall temperature} thì không, dù không luật ngữ pháp nào cấm? Vì sao \emph{make a decision} đúng còn \emph{do a decision} thì không? Đây không phải chuyện nghĩa, cũng không phải chuyện ngữ pháp --- nó là một lớp riêng, và nó là thứ phân biệt rõ nhất người viết thành thạo với người viết đúng luật. Chương \ptr{D/17} nói về lớp đó.
\ifSubfilesClassLoaded{\printbibliography[title={Nguồn}]}{}
\end{document}
